\part{Hypergeometric Functions}

\section{Gauss Hypergeometric Functions}

\begin{exercise}[Asymptotic translation symmetry and scattering channels]
Let $\R$ act on its configuration line by translations and on $L^2(\R)$ by
pullback.  Let
\[
 H_0=-\frac{\hbar^2}{2m}\frac{\dd^2}{\dd x^2},\qquad H=H_0+V,
\]
where $V$ is real and
$\int_{\R}(1+\abs x)\abs{V(x)}\dd x<\infty$.  Construct $H$ from its closed
semibounded quadratic form after defining densely defined, closed, and
semibounded forms and proving their self-adjoint representation theorem.
For a semibounded symmetric operator, close its form and construct its
Friedrichs extension.  Prove that $H_0$ is translation-invariant and
that $V$ breaks this symmetry only by a short-range interaction, and prove
that the absolutely continuous spectrum
on $(0,\infty)$ has two asymptotic translation channels.  Construct the wave
operators
\[
 \Omega_\pm=\operatorname*{s-lim}_{t\to\pm\infty}
 e^{\ii Ht/\hbar}e^{-\ii H_0t/\hbar}
\]
first on Fourier transforms of compactly supported functions away from zero,
using the Duhamel formula and the one-dimensional dispersive estimate, and
then by density.  Prove isometry, intertwining, and completeness from the
Volterra resolvent.  Define $S=\Omega_+^*\Omega_-$ and prove that it commutes
with $H_0$.  Show that time reversal makes the
two-channel scattering matrix symmetric and that an even potential further
decomposes it into parity channels.
\end{exercise}

\begin{exercise}[Jost solutions and the scattering matrix]
For $k>0$ define the Jost solutions as the unique solutions of
\[
 -\frac{\hbar^2}{2m}f''+Vf=\frac{\hbar^2k^2}{2m}f,
 \qquad f_\pm(k,x)e^{\mp\ii kx}\longrightarrow1
 \quad(x\to\pm\infty),
\]
and construct them from their Volterra integral equations.  Use their
Wronskians to define transmission and left and right reflection amplitudes,
fixing the convention
\[
 f_+(k,x)=t(k)^{-1}e^{\ii kx}
 +t(k)^{-1}r_-(k)e^{-\ii kx}+o(1)
 \quad(x\to-\infty).
\]
Prove that these coefficients form the fiber of $S$, derive unitarity and
reciprocity, and meromorphically continue them to $\Im k>0$.  Prove that poles
of $t(k)$ at $k=\ii\kappa$ are bound states of energy
$-\hbar^2\kappa^2/(2m)$, derive their norming constants from the residues, and
identify lower-half-plane continuation poles as resonances.
\end{exercise}

\begin{exercise}[P\"oschl--Teller hypergeometric functions]
For
\[
 V(x)=-\frac{\hbar^2}{2ma^2}\lambda(\lambda+1)
 \operatorname{sech}^2(x/a),\qquad \lambda>0,
\]
put $z=(1-\tanh(x/a))/2$ and reduce the Jost equation to the Gauss equation.
Identify its Gauss hypergeometric parameters and specialize the Euler integral
and the connection formula between $z=0$ and $z=1$.
Express both Jost solutions through these functions and obtain, with $p=ka$,
\[
 t(p)=\frac{\Gamma(1+\lambda-\ii p)\Gamma(-\lambda-\ii p)}
 {\Gamma(1-\ii p)\Gamma(-\ii p)},
\]
in the preceding asymptotic convention.  Use the gamma
reflection formula to prove
\[
 \abs{t(p)}^2=
 \frac{\sinh^2(\pi p)}{\sinh^2(\pi p)+\sin^2(\pi\lambda)},\qquad
 \abs{r(p)}^2=
 \frac{\sin^2(\pi\lambda)}{\sinh^2(\pi p)+\sin^2(\pi\lambda)}.
\]
\end{exercise}

\begin{exercise}[Reflectionless transmission spectroscopy]
For integral $\lambda=N$, prove that the P\"oschl--Teller potential is
reflectionless at every positive energy and that its bound-state poles give
\[
 E_n=-\frac{\hbar^2}{2ma^2}(N-n)^2,\qquad n=0,\ldots,N-1.
\]
For nonintegral $\lambda$, locate the bound-state poles with
$n\in\Z_{\ge0}$ and $n<\lambda$ and the continued resonance poles, and derive the threshold
and high-energy asymptotics of the reflection probability.  Formulate a
one-dimensional matter-wave or microwave-guide experiment that measures
transmitted and reflected flux as a function of $k$: the model predicts the
displayed transmission curve, exact reflectionlessness at integral $N$, and
the bound-state line positions.  Separate calibration of $a$ and the potential
depth, finite preparation width, absorption, and truncation of the ideal
$\operatorname{sech}^2$ profile from these predicted observables.
\end{exercise}

\section{Legendre Functions}

\begin{exercise}[Hyperbolic isometries and the Casimir]
Let $SL_2(\R)$ act on $\mathbb H=\{x+\ii y:y>0\}$ by M\"obius maps.  Prove
that it preserves $\dd s^2=(\dd x^2+\dd y^2)/y^2$, identify the stabilizer of
$\ii$ with $SO(2)$, and derive the Iwasawa and Cartan decompositions.  With
the nonnegative convention define
\[
 \Delta_{\mathbb H}=-y^2(\partial_x^2+\partial_y^2)
\]
and show that it is the action of the quadratic Casimir on right-$SO(2)$
invariant functions.
\end{exercise}

\begin{exercise}[Principal, complementary, and discrete series]
Construct the spherical principal series by unitary induction from the upper
triangular subgroup, the complementary series by its invariant kernel, and
the holomorphic and antiholomorphic discrete series on weighted Bergman
spaces.  Determine their Casimir and $SO(2)$ spectra, including the signs of
the discrete-series weights.  Prove irreducibility by cyclicity of a
$K$-finite vector and classify the irreducible unitary representations by
the possible lowest/highest weight or continuous Casimir parameter.
\end{exercise}

\begin{exercise}[$SL_2(\R)$ Plancherel decomposition]
Construct the standard intertwining operator between opposite principal
series, compute its scalar on every $SO(2)$ type from the raising and lowering
relations, and determine its poles and residues.  Shift its contour to recover
the discrete-series summands and derive the group Plancherel measure, with
both even and odd principal-series densities.  Prove inversion first for
compactly supported smooth $K$-finite functions and extend it to $L^2$.
\end{exercise}

\begin{exercise}[Legendre functions and the Mehler--Fock transform]
With fixed branches, use the Gauss hypergeometric function to define
\[
 P_\nu^\mu(z)=\frac1{\Gamma(1-\mu)}
 \left(\frac{z+1}{z-1}\right)^{\mu/2}
 {}_2F_1\!\left(-\nu,\nu+1;1-\mu;\frac{1-z}{2}\right).
\]
Derive the associated Legendre equation and show that the spherical function
at spectral parameter $\lambda$ is
$P_{-1/2+\ii\lambda}(\cosh r)$.  From the radial Sturm--Liouville problem
prove the Mehler--Fock inversion and Plancherel formulas, including the
density $\lambda\tanh(\pi\lambda)$.  Derive the analogous fixed-$SO(2)$-type
transform with $P_{-1/2+\ii\lambda}^{-m}$.
\end{exercise}

\begin{exercise}[Hyperbolic-drum spectra]
Let a membrane occupy a geodesic disk of curvature radius $a$ and hyperbolic
radius $R$, with wave speed $v$ and Dirichlet boundary.  Separate its wave
equation into $SO(2)$ types and prove that its resonances are
\[
 \omega_{m,k}=\frac va\sqrt{\frac14+\lambda_{m,k}^2},
 \qquad
 P_{-1/2+\ii\lambda_{m,k}}^{-\abs m}(\cosh R)=0.
\]
Prove the $m\ne0$ double degeneracy and the Weyl asymptotic from the
Mehler--Fock spectral measure.  The zero locations, degeneracies, and area
coefficient are the three measurable predictions.
\end{exercise}

\section{Kummer Functions}

\begin{exercise}[Runge--Lenz symmetry of the Coulomb Hamiltonian]
Let $V$ be an oriented three-dimensional Euclidean space, let $X$ and $P$ be
the position and momentum operators in the Schr\"odinger representation, and
use the metric and orientation to define their cross product.  Construct the
attractive Coulomb Hamiltonian by its closed semibounded form,
\[
 H=\frac{P^2}{2m}-\frac\kappa{\abs X},\qquad \kappa>0,
\]
and on a common invariant core define
\[
 L=X\times P,
 \qquad
 A=\frac1{2m}(P\times L-L\times P)-\kappa\frac X{\abs X}.
\]
Prove the Hardy inequality needed to make the Coulomb term form-bounded with
relative bound zero and hence justify the stated self-adjoint realization.
For $a,b\in V$, derive
\[
 [a\cdot L,b\cdot L]=\ii\hbar(a\times b)\cdot L,
 \quad [a\cdot L,b\cdot A]=\ii\hbar(a\times b)\cdot A,
 \quad [a\cdot A,b\cdot A]
 =-\frac{2\ii\hbar H}{m}(a\times b)\cdot L,
\]
together with $L\cdot A=0$ and
$A^2=\kappa^2+2H(L^2+\hbar^2)/m$.  On the negative spectral subspace rescale
$A$ and prove that $L$ and the rescaled Runge--Lenz operator generate
$\mathfrak{so}(4)\simeq\mathfrak{su}(2)\oplus\mathfrak{su}(2)$.  Use its two
Casimirs to derive
\[
 E_n=-\frac{m\kappa^2}{2\hbar^2n^2},
 \qquad \operatorname{mult}(E_n)=n^2,qquad n\geq1.
\]
\end{exercise}

\begin{exercise}[Kummer and Laguerre spectral functions]
Specialize the generalized hypergeometric series and verify its Gauss
confluence formula
\[
 {}_1F_1(a;c;z)=\lim_{b\to\infty}{}_2F_1(a,b;c;z/b)
 =\sum_{j\geq0}\frac{(a)_j}{(c)_j j!}z^j.
\]
Derive Kummer's equation, its Euler integral and connection formula, and the
confluence of the Gauss singular points, initially under
$\Re c>\Re a>0$ for the integral and then by meromorphic continuation.  For
$\alpha>-1$ define
\[
 L_j^\alpha(z)=\frac{(\alpha+1)_j}{j!}
 {}_1F_1(-j;\alpha+1;z),
\]
and derive its Rodrigues formula, recurrence, differential equation,
orthogonality for $z^\alpha e^{-z}\dd z$, squared norm, and completeness.

Put $a_0=\hbar^2/(m\kappa)$.  Separate the Coulomb Hamiltonian by the
spherical decomposition and prove that its normalized bound radial functions
are
\[
 R_{n\ell}(r)=N_{n\ell}e^{-r/(na_0)}
 \left(\frac{2r}{na_0}\right)^\ell
 L_{n-\ell-1}^{2\ell+1}\!\left(\frac{2r}{na_0}\right),
 \quad 0\leq\ell<n.
\]
Determine $N_{n\ell}$, prove orthogonality and completeness on the bound
spectral subspace, and identify the resulting $SO(4)$ decomposition with the
tensor product of the two spin-$(n-1)/2$ representations obtained from the
Runge--Lenz symmetry.
\end{exercise}

\begin{exercise}[Coulomb scattering functions]
For the repulsive potential $V(r)=\kappa/r$, set $\rho=kr$ and
$\eta=m\kappa/(\hbar^2k)$.  Define the regular Coulomb function by
\[
 F_\ell(\eta,\rho)=C_\ell(\eta)\rho^{\ell+1}e^{-\ii\rho}
 {}_1F_1(\ell+1-\ii\eta;2\ell+2;2\ii\rho),
\]
choosing $C_\ell(\eta)>0$ so that its large-$\rho$ sine amplitude is one.
Derive the radial equation and specialize the Kummer connection formula to
obtain
$\sigma_\ell=\arg\Gamma(\ell+1+\ii\eta)$ and
$\sigma_\ell-\sigma_{\ell-1}=\arctan(\eta/\ell)$.  Sum the resulting
Legendre series by its generating function and prove that its long-range
phase is compatible with flux conservation.
\end{exercise}

\begin{exercise}[Hydrogen spectroscopy and Rutherford scattering]
For the attractive potential with
$\kappa=Ze^2/(4\pi\epsilon_0)$, use the Laguerre spectral functions and the
Wigner--Eckart theorem to derive the electric-dipole rules
$\Delta\ell=\pm1$ and $\Delta m=0,\pm1$, and express every radial line
strength by a terminating Laguerre integral.  Prove that the line frequencies
are
\[
 \nu_{n\to n'}=\frac{m\kappa^2}{4\pi\hbar^3}
 \left|\frac1{n^2}-\frac1{{n'}^2}\right|,
\]
with the reduced mass replacing $m$ for a moving nucleus.

For the repulsive potential, use the Coulomb connection coefficient to derive
the amplitude, up to an unobservable constant phase,
\[
 f_C(\theta)=-\frac{\eta}{2k\sin^2(\theta/2)}
 \exp\!\left[-2\ii\eta\log\sin\frac\theta2+2\ii\sigma_0\right],
\]
and hence
\[
 \frac{\dd\sigma}{\dd\Omega}=
 \left(\frac{Z_1Z_2e^2}{16\pi\epsilon_0E}\right)^2
 \csc^4\frac\theta2.
\]
Identify spectroscopic frequencies and intensities, beam energy, scattering
angle, and detector solid angle as measured quantities.  State separately
the nonrelativistic, point-nucleus, spin-free, and radiative limitations of
the two predictions.
\end{exercise}

\section{Heun Functions}

\begin{exercise}[Lorentzian geometry and spherical vacuum metrics]
Define a Lorentzian metric, its Levi--Civita connection by the Koszul identity,
curvature, Ricci tensor, and geodesics.  Derive the Jacobi equation from a
variation of geodesics.  For a static spherically symmetric metric, use the
area radius and proper-time normalization at infinity to solve
$\operatorname{Ric}=0$, obtaining
\[
 g=-\left(1-\frac{r_s}{r}\right)c^2\dd t^2
 +\left(1-\frac{r_s}{r}\right)^{-1}\dd r^2+r^2\dd\Omega^2,
 \qquad r_s=\frac{2GM}{c^2}.
\]
Prove that its connected exterior isometry group is
$\R\times SO(3)$.
\end{exercise}

\begin{exercise}[Regge--Wheeler decomposition]
Set $G=c=1$ temporarily and define the tortoise coordinate by
$\dd r_*/\dd r=(1-2M/r)^{-1}$.  Separate the scalar wave equation using the
$SO(3)$ decomposition and prove that $u_{\ell m}=r\phi_{\ell m}$ satisfies
\[
 \left(-\partial_t^2+\partial_{r_*}^2-V_\ell(r)\right)u_{\ell m}=0,
 \qquad
 V_\ell(r)=\left(1-\frac{2M}{r}\right)
 \left(\frac{\ell(\ell+1)}{r^2}+\frac{2M}{r^3}\right).
\]
Construct the conserved Wronskian and the ingoing/outgoing scattering
coefficients, and prove their flux identity.
\end{exercise}

\begin{exercise}[Confluent Heun functions and quasinormal modes]
Define a confluent Heun function as the Frobenius solution at a regular
singular point of
\[
 y''+\left(\alpha+\frac{\beta+1}{z}+\frac{\gamma+1}{z-1}\right)y'
 +\left(\frac\mu z+\frac\nu{z-1}\right)y=0,
\]
with its parameters and initial exponent specified.  Transform the separated
Schwarzschild radial equation to this form and derive the three-term
recurrence of its Frobenius coefficients.  With time dependence
$e^{-\ii\omega t}$, define a quasinormal frequency as a zero of the
connection coefficient multiplying the incoming solution at
infinity for the solution ingoing at the horizon.  Express that coefficient
as the limit of the recurrence's continued fraction and prove its convergence
under the minimal-solution condition.
\end{exercise}

\begin{exercise}[Schwarzschild scalar ringdown]
For $\ell\to\infty$, locate the maximum of $V_\ell$ at the unstable photon
orbit $r=3M$.  Define the parabolic-cylinder function $D_\nu(z)$ as the
solution of $y''+(\nu+\tfrac12-z^2/4)y=0$ normalized by
$D_\nu(z)\sim z^\nu e^{-z^2/4}$ for $\abs{\arg z}<3\pi/4$.  Derive its
connection formula by analytically continuing
\[
 D_\nu(z)=\frac{e^{-z^2/4}}{\Gamma(-\nu)}
 \int_0^\infty e^{-zt-t^2/2}t^{-\nu-1}\dd t,
 \qquad \Re\nu<0,
\]
and match these functions
across the potential maximum to obtain, for fixed overtone $n$,
\[
 \frac{GM\omega_{\ell n}}{c^3}=
 \frac{\ell+1/2-\ii(n+1/2)}{3\sqrt3}+O(\ell^{-1}).
\]
For the scalar $\ell=2$ fundamental mode, truncate the confluent-Heun
continued fraction with a rigorously bounded minimal-solution tail and obtain
\[
 \frac{GM\omega_{20}}{c^3}
 =0.48364\ldots-0.09676\ldots\ii.
\]
Show that the predicted scalar waveform has oscillation frequency
$\Re\omega/(2\pi)$ and damping time $1/\abs{\Im\omega}$.  Explain how the
continued-fraction value supplies the finite-$\ell$ correction and compare it
with the eikonal approximation.  State that this is a test-field prediction,
not the gravitational $\ell=2$ mode, and identify the mass calibration and
measured oscillation and damping times.
\end{exercise}

\part{Character Theory}

\section{Weyl Characters}

\begin{exercise}[Lie algebras, enveloping algebras, and Hopf algebras]
Define a complex Lie algebra, Lie ideals, representations, and the adjoint
representation.  Construct the tensor algebra by its universal property and
define
\[
 U(\mathfrak g)=T(\mathfrak g)/
 \langle x\otimes y-y\otimes x-[x,y]\rangle
\]
by its universal property among associative algebras receiving a Lie map from
$\mathfrak g$.  Filter $U(\mathfrak g)$ by tensor degree, construct the
canonical graded map
$\operatorname{Sym}(\mathfrak g)\to\operatorname{gr}U(\mathfrak g)$, and
prove it is an isomorphism by the rewriting filtration and the Jacobi
identity.  Deduce the Poincar\'e--Birkhoff--Witt theorem, induction and
restriction of modules.  Define a triangular decomposition
$\mathfrak g=\mathfrak n_-\oplus\mathfrak h\oplus\mathfrak n_+$ by requiring
$\mathfrak h$ to be abelian,
$[\mathfrak h,\mathfrak n_\pm]\subseteq\mathfrak n_\pm$, and
$\mathfrak b_\pm=\mathfrak h\oplus\mathfrak n_\pm$ to be subalgebras;
construct the associated Verma modules by induction.  Define coalgebras,
bialgebras, and Hopf algebras by
their universal diagrams, and construct the cocommutative Hopf structure on
$U(\mathfrak g)$ for which every $x\in\mathfrak g$ is primitive.
\end{exercise}

\begin{exercise}[Root data and Weyl groups]
Define an affine complex algebraic group by its finitely generated reduced
commutative coordinate Hopf algebra, after defining nilpotent elements and
reduced algebras.  Define connectedness, unipotent groups,
the unipotent radical, and reductivity, and construct its Lie algebra from
derivations at the counit.  For a connected complex reductive group $G$,
define a maximal torus $T$, its
character and cocharacter lattices, roots, coroots, and the Weyl group
$W=N_G(T)/T$.  Prove that $W$ is finite, construct a $W$-invariant Euclidean
metric on the real weight space $X^*(T)\otimes_\Z\R$ by averaging, and prove
that every root reflection is the orthogonal reflection in its root
hyperplane.  Derive the root-space decomposition of $\mathfrak g$ and the
rank-one $\mathfrak{sl}_2$ subalgebra associated with each root.  Choose a
positive system, define simple roots, dominant weights, the Weyl length, and
$\rho$, and prove the chamber and Bruhat decompositions.
\end{exercise}

\begin{exercise}[Highest-weight decomposition]
Define weight spaces without choosing vectors in them.  Prove that every
finite-dimensional irreducible $G$-representation has a unique dominant
highest weight.  For a Borel algebra
$\mathfrak b=\mathfrak t\oplus\mathfrak n_+$ define
$M_\lambda=U(\mathfrak g)\otimes_{U(\mathfrak b)}\C_\lambda$ and construct
the irreducible representation as its unique irreducible quotient.  Use the
rank-one strings to prove finite
dimensionality exactly for dominant integral weights, complete reducibility,
and the highest-weight classification.
\end{exercise}

\begin{exercise}[Weyl characters and Schur functions]
Define exterior powers as alternating tensor quotients and define the
determinant of an endomorphism as the scalar by which it acts on the highest
nonzero exterior power.  Define a partition as a finite weakly decreasing
sequence of nonnegative integers.  In the group ring of the weight lattice
define
$A_\mu=\sum_{w\in W}(-1)^{\ell(w)}e^{w\mu}$.  Construct the BGG resolution
from Verma modules and use its Euler character to prove
\[
 \operatorname{ch}V_\lambda=\frac{A_{\lambda+\rho}}{A_\rho},
 \qquad
 A_\rho=e^\rho\prod_{\alpha>0}(1-e^{-\alpha}).
\]
For $G=GL_r$ identify these characters with the Schur functions
\[
 s_\lambda(x)=
\frac{\det(x_i^{\lambda_j+r-j})}{\det(x_i^{r-j})},
\]
and derive the tableaux expansion, Cauchy identity, and
Littlewood--Richardson tensor-product rule.
\end{exercise}

\begin{exercise}[Flavor multiplets and the Gell-Mann--Okubo prediction]
Use the $SU(3)$ Weyl character $s_{(2,1)}$ to recover the eight weights of the
adjoint representation with their multiplicities.  Let isospin $SU(2)$ and
hypercharge $U(1)$ be the evident subgroup and show that the octet decomposes
as two isodoublets, one isotriplet, and one isosinglet.  Assume the leading
symmetry-breaking mass operator is an isoscalar in $1\oplus8$.  Apply the
Wigner--Eckart theorem and eliminate its three reduced parameters to derive
\[
 2(M_N+M_\Xi)=3M_\Lambda+M_\Sigma.
\]
Explain which four observed baryon centroid masses test the relation and
separate the exact representation-theoretic multiplicities from the
first-order symmetry-breaking assumption.
\end{exercise}

\section{Automorphic Forms}

\begin{exercise}[Modular group and its fundamental domain]
Define $SL_2(\Z)$, its action on $\mathbb H$, and
$S\tau=-1/\tau$, $T\tau=\tau+1$.  Prove the presentation
$PSL_2(\Z)=\langle S,T\mid S^2=(ST)^3=1\rangle$ by reducing points to the
standard fundamental domain.  Compute its hyperbolic area and the cusp,
elliptic-point, and orbifold stabilizer data.
\end{exercise}

\begin{exercise}[Eisenstein series and the scattering coefficient]
Construct $L^2(SL_2(\Z)\backslash\mathbb H)$ with hyperbolic measure and the
nonnegative Laplacian.  Define the Eisenstein series
$E(\tau,s)=\sum_{\Gamma_\infty\backslash\Gamma}(\Im\gamma\tau)^s$, where
$\Gamma_\infty$ is the stabilizer of the cusp at infinity, for $\Re s>1$.
Define the zeta function $\zeta(s)=\sum_{n\geq1}n^{-s}$ there and its
completion $\xi(s)=\pi^{-s/2}\Gamma(s/2)\zeta(s)$.  Unfold the Eisenstein
series, continue it using the theta Mellin transform, and prove that its
constant term is
\[
 E(x+\ii y,s)_0=y^s+\varphi(s)y^{1-s},\qquad
 \varphi(s)=\frac{\xi(2s-1)}{\xi(2s)}.
\]
Derive its functional equation and
$\varphi(s)\varphi(1-s)=1$, including unitarity on $\Re s=1/2$.
\end{exercise}

\begin{exercise}[Automorphic spectral and trace decomposition]
Develop the Friedrichs realization of the nonnegative Laplacian at the cusp.
Truncate the Eisenstein series, prove the Maass--Selberg relation, and obtain
the orthogonal decomposition into constants, cusp forms, and the continuous
Eisenstein integral.  Build the automorphic kernel from a compactly supported
point-pair invariant, unfold it by conjugacy classes, and derive the Selberg
trace identity with its identity, elliptic, parabolic, and hyperbolic terms.
Recover the Weyl area term.
\end{exercise}

\begin{exercise}[Modular forms, theta constants, and the eta function]
Define modular forms by their weight, multiplier, and cusp growth.  Starting
from the theta transformation proved in Fourier Analysis, define
$\vartheta_2,\vartheta_3,\vartheta_4$ by their shifted or signed lattice sums
and derive all their $S,T$ transformations without re-proving Poisson
summation.  Define
$\eta(\tau)=q^{1/24}\prod_{n\ge1}(1-q^n)$ with
$q=e^{2\pi\ii\tau}$, and derive its transformations from the Jacobi triple
product.  Construct $E_4,E_6$, prove the valence formula, and show
\[
 \C[E_4,E_6]=\bigoplus_{k\ge0}M_k,\qquad
 \Delta=\eta^{24}=\frac{E_4^3-E_6^2}{1728}.
\]
\end{exercise}

\begin{exercise}[Eisenstein scattering on the modular surface]
Model a scalar waveguide by the modular surface with its cusp attached to a
single asymptotic lead.  At spectral value
$s=\tfrac12+\ii r$ with $r>0$, so that the nonnegative Laplacian has value
$\tfrac14+r^2$, identify $y^{1/2-\ii r}$ and $y^{1/2+\ii r}$ as the two
cusp waves.  Orient the lead so that the $y^s$ coefficient is the unit
incoming amplitude and the $y^{1-s}$ coefficient is outgoing, and prove that
the reflection amplitude is
$\varphi(\tfrac12+\ii r)$.  Write
$\varphi(\tfrac12+\ii r)=e^{2\ii\delta(r)}$ and, for wave speed $v$ and
$\omega(r)=v\sqrt{\tfrac14+r^2}$, derive the Wigner delay
\[
 \tau_W(r)=2\frac{\dd\delta}{\dd\omega}
 =\frac{2\sqrt{\tfrac14+r^2}}{v r}\frac{\dd\delta}{\dd r}.
\]
Show that the meromorphically continued poles of $\varphi$ are the scattering
resonances.  Thus the completed-zeta quotient predicts the reflected phase,
time delay, and resonance locations; separate the constant-curvature,
single-cusp, scalar-wave, and lead-matching assumptions from those measured
quantities.
\end{exercise}

\section{Virasoro Characters}

\begin{exercise}[The Virasoro central extension]
Let $\operatorname{Diff}^+(S^1)$ act on the circle and on scalar fields by
pullback, and identify its infinitesimal action with smooth vector fields.
After complexification on the punctured complex line let
$L_n=-z^{n+1}\partial_z$, $n\in\Z$, be the distinguished Fourier vector
fields.  Prove
$[L_m,L_n]=(m-n)L_{m+n}$.  Define
\[
 \omega(L_m,L_n)=\frac{m^3-m}{12}\delta_{m+n,0}
\]
and prove directly that it is alternating and satisfies the Lie-algebra
cocycle identity.  Define
$\operatorname{Vir}=\bigoplus_{n\in\Z}\C L_n\oplus\C C$ by
\[
 [L_m,L_n]=(m-n)L_{m+n}+\omega(L_m,L_n)C,
 \qquad [C,L_n]=0,
\]
and verify the involution $L_n^*=L_{-n}$, $C^*=C$.  Recover the cocycle as
\[
 \omega(f\partial_z,g\partial_z)
 =\frac1{12}\operatorname{Res}_{z=0}(f'''(z)g(z)\dd z),
\]
compute its change under a new circle coordinate as a coboundary, and deduce
that the central-extension class is coordinate-independent.
\end{exercise}

\begin{exercise}[Positive-energy highest-weight decomposition]
Call a unitary Virasoro representation positive-energy when $C$ acts by a
real scalar $c$, $L_0$ has spectrum bounded below with finite-dimensional
eigenspaces, and the algebraic sum of those eigenspaces is dense.  Use
$[L_0,L_n]=-nL_n$ and the involution to prove that every irreducible such
representation is generated by a vector $v$ satisfying
\[
 L_0v=hv,\qquad L_nv=0\ (n>0),\qquad Cv=cv,
\]
and prove the orthogonal decomposition of an arbitrary positive-energy
representation into irreducible highest-weight summands by induction over
its energy levels.  Define the Verma module $M(c,h)$ by this universal
highest-weight property.  Apply the PBW theorem to show that its level-$N$
subspace is indexed by partitions of $N$; write $p(N)$ for their number and
put $p(N)=0$ for $N<0$.  Derive
\[
 \operatorname{ch}M(c,h)
 =\operatorname{Tr}_{M(c,h)}q^{L_0-c/24}
 =\frac{q^{h-(c-1)/24}}{\eta(\tau)},
 \qquad q=e^{2\pi\ii\tau}.
\]
\end{exercise}

\begin{exercise}[Shapovalov determinants and unitary representations]
Give $M(c,h)$ the contravariant form determined by
$\langle v,v\rangle=1$ and $L_n^*=L_{-n}$.  Compute its Gram determinants
at levels one and two.  Put
\[
 c=1-6(b-b^{-1})^2,
 \qquad
 h_{r,s}(b)=\frac{(r/b-sb)^2-(b^{-1}-b)^2}{4}.
\]
Using the PBW filtration, construct a singular vector when
$h=h_{r,s}(b)$, compare the degree and leading coefficient of the level-$N$
determinant as a polynomial in $h$, and prove
\[
 \det_N(c,h)=C_N
 \prod_{\substack{r,s\geq1\\rs\leq N}}
 (h-h_{r,s}(b))^{p(N-rs)}
\]
with $C_N>0$ fixed by the ordered PBW monomials.  Track the signs between
consecutive zero curves and quotient the null spaces to prove
that irreducible positive-energy modules are unitary for $c\geq1$, $h\geq0$,
and, below $c=1$, precisely at
\[
 c_m=1-\frac6{m(m+1)},\qquad
 h_{r,s}^{(m)}=\frac{((m+1)r-ms)^2-1}{4m(m+1)}
\]
with $m\geq2$, $1\leq r\leq m-1$, $1\leq s\leq m$, modulo
$(r,s)\sim(m-r,m+1-s)$.
\end{exercise}

\begin{exercise}[Minimal-model characters]
For $K\in\N$ and $\mu\in\Z$ define the level theta series
\[
 \Theta_{\mu,K}(\tau)
 =\sum_{n\in\Z}q^{(2Kn+\mu)^2/(4K)}.
\]
Determine the inclusions among the singular Verma submodules at
$(c_m,h_{r,s}^{(m)})$, resolve the irreducible quotient by their alternating
chain, and take its Euler character to prove
\[
 \chi_{r,s}^{(m)}(\tau)=
 \frac{\Theta_{(m+1)r-ms,\,m(m+1)}(\tau)
 -\Theta_{(m+1)r+ms,\,m(m+1)}(\tau)}{\eta(\tau)}.
\]
Use the Poisson summation formula to derive the finite
$S$- and $T$-matrices of these characters.  At $m=3$ identify the three
weights $0,1/2,1/16$ and prove
\[
 \chi_0=\frac12\left(\sqrt{\frac{\vartheta_3}{\eta}}
 +\sqrt{\frac{\vartheta_4}{\eta}}\right),\quad
 \chi_{1/2}=\frac12\left(\sqrt{\frac{\vartheta_3}{\eta}}
 -\sqrt{\frac{\vartheta_4}{\eta}}\right),\quad
 \chi_{1/16}=\frac1{\sqrt2}\sqrt{\frac{\vartheta_2}{\eta}},
\]
fixing the square roots by their positive $q$-expansions.
\end{exercise}

\begin{exercise}[Virasoro characters and monstrous moonshine]
Define a vertex operator algebra as a positive-energy Virasoro representation
with vacuum $\mathbf1$, translation operator $D$, conformal vector, and
state--field map
$Y(a,z)=\sum_{n\in\Z}a_{(n)}z^{-n-1}$ satisfying
\[
 Y(\mathbf1,z)=1,\qquad
 Y(a,z)\mathbf1\in a+zV[[z]],\qquad
 [D,Y(a,z)]=\partial_zY(a,z),
\]
and $(z-w)^N[Y(a,z),Y(b,w)]=0$ for sufficiently large $N$, with the
conformal field equal to $T(z)$.  Derive the mode Jacobi identity and prove
that the graded dimension is a Virasoro character.  Define $j=E_4^3/\Delta$
and use the
valence formula to show that every modular function holomorphic away from the
cusp is a polynomial in $j$.  Compute
\[
 J(\tau)=j(\tau)-744=q^{-1}+196884q+21493760q^2+\cdots.
\]
Write $J(q)=\sum_{n\geq-1}c(n)q^n$.  Construct the graded traces of the
moonshine vertex operator algebra and prove its Borcherds denominator identity
\[
 p^{-1}\prod_{m>0,\,n\in\Z}(1-p^mq^n)^{c(mn)}=J(p)-J(q).
\]
Extract the denominator recursion and use the vacuum normalization and first
graded pieces to prove that its identity character is $J$.
Verify
$196884=1+196883$ and
$21493760=1+196883+21296876$ as decompositions into irreducible Monster
dimensions.
\end{exercise}

\begin{exercise}[Conformal spectra and the Cardy prediction]
Define a unitary two-dimensional conformal theory on a circle of circumference
$L$ by commuting positive-energy Virasoro actions with central charges
$c,\bar c$ and weights $h,\bar h$, and use units $\hbar=1$ in conformal
finite-size formulas.  Define $T(z)=\sum_nL_nz^{-n-2}$ and transform it from
the plane to the cylinder, deriving the Schwarzian contribution from the
Virasoro cocycle.  Obtain the Hamiltonian and momentum
\[
 H=\frac{2\pi v}{L}\left(L_0+\bar L_0-\frac{c+\bar c}{24}\right),
 \qquad P=\frac{2\pi}{L}(L_0-\bar L_0).
\]
For a finite character family and nonnegative integral multiplicity matrix
$M$, prove that the torus trace
$Z=\sum_{i,j}M_{ij}\chi_i\overline{\chi_j}$ is modular invariant exactly
when $M$ commutes with its $S,T$ matrices; identify the integral-spin and
$c-\bar c\in24\Z$ conditions imposed by $T$.  If $h_{\min}$ is the least
chiral weight coupled to the vacuum, define
$c_{\rm eff}=c-24h_{\min}$.  Apply the $S$ transformation and a saddle-point
estimate to prove
\[
 \log\rho(\Delta)=2\pi\sqrt{\frac{c_{\rm eff}\Delta}{6}}+o(\sqrt\Delta).
\]
Thus predict both the finite-size ground-energy shift
$-\pi v(c+\bar c)/(12L)$ and the high-energy level-counting rate, specifying
$v,L$, and the measured energies.
\end{exercise}
